\documentclass[11pt,a4paper]{article}
\usepackage[margin=25mm]{geometry}
\usepackage{amsmath,amssymb,amsthm,hyperref}
\setlength{\emergencystretch}{2em}
\newtheorem{lemma}{Lemma}
\newtheorem{theorem}[lemma]{Theorem}
\newcommand{\gen}[1]{\langle#1\rangle}
\title{An alternating group on the pairs}
\author{Human Compression project: a component of the L044 recovery}
\date{}
\begin{document}\maketitle
Let $G=A_{2m}=\gen{\rho_0,\ldots,\rho_{r-1}}$ in its
natural action, with a string C-marking. Suppose a maximal
omission $H=G_i$ is transitive and preserves $m$ pairs, and
its induced block group is $A_m$. Throughout $m\ge8$.
We assume explicitly the rank bound for natural alternating
groups of every smaller degree $12\le q<2m$. Degrees below
twelve use the acquired complete finite maxima. These are the
strong-induction premises, not the theorem at degree $2m$.

We acquire the complete L047 block decomposition, its exact
transitive-rank theorem and finite exceptions, the extremal
independent-string theorem for symmetric groups, and the
accepted lemma on centralizers of primitive actions on pairs.
In particular, in a putative case $r\ge m$ the minimal
block-generating family $L$ is a connected interval, its
commuting remainder $C$ has size at most one, and its other
remainder $R$ has size at most one. The precise actual-action
proof of this last fact is retained.

Write $V=\mathbb F_2^m$, $z=(1,\ldots,1)$, and let $E$
be the even-weight subspace. The acquired elementary module
argument says that the $A_m$-invariant subspaces of $V$ are
$0,\gen z,E,V$ for $m\ge5$: a three-cycle applied to a
nonconstant vector gives a difference of weight two, and
conjugating such a difference generates $E$. The kernel of
an even permutation group on the pairs lies in $E$.

\begin{lemma}\label{parityproduct}
Let $m\ge8$ be even. Let $Q\le A_m$ be intransitive and
have a string C-marking of rank $m-3$. Suppose this marking
arises faithfully from a permutation group $T$ on $2m$
points preserving $m$ pairs, with block image $Q$.
Assume the stated smaller-degree alternating rank bounds.
Then every even involution centralizing $T$ preserves the
pairs.
\end{lemma}
\begin{proof}
For an intransitive group on orbit sizes $a_1,\ldots,a_q$,
successive projection and kernel normalization give independent
length at most $\sum(a_j-1)=m-q$, by Whiston's bound.
Equality forces the full product of the symmetric groups on
those orbits. Indeed, choose a minimal subfamily generating
one projection and normalize the remaining members into its
kernel by words in that subfamily. Equality forces both
bounds to be attained; Whiston's equality clause identifies
the first projection and the last kernel as the full symmetric
factors. Iterate over the remaining orbits. This also handles
singleton orbits, whose contribution is zero.

Consequently $Q$ has at most three orbits. Three would force
that full product, containing an odd permutation, contrary
to $Q\le A_m$. Thus there are exactly two orbits, of sizes
$a\le b$, with $a+b=m$.
If $a=1$, the faithful action on the other $b=m-1$ points
has rank $b-2$. Here $b\ge7$ is odd. The complete acquired
L047 theorem excludes every proper transitive group other
than $A_b,S_b$ at this rank; its proper high-rank exceptions
have even degree. The even condition excludes $S_b$, and
the assumed or finite alternating bound excludes $A_b$.
Hence $a\ge2$.

Projecting onto the $a$ points and normalizing the other
generators into the kernel gives
\[
 m-3\le\mu(Q|_a)+\mu'(A_b)\le(a-1)+(b-2)=m-3.
\]
Here $\mu$ is maximum irredundant generating length and
$\mu'$ maximum independent length. Equality makes the
first projection $S_a$. Reversing the two projections makes
the other $S_b$. Thus $Q$ is subdirect in $S_a\times S_b$.
If $b\ge5$, its kernel in the second factor is a normal
subgroup of $S_b$ contained in $A_b$, hence is $1$ or $A_b$.
The first possibility would identify $Q$ with $S_a$, of
independent length at most $a-1<m-3$. Therefore that kernel
is $A_b$. Projection onto the first factor and the parity
condition now identify the full group:
\[
                     Q=(S_a\times S_b)\cap A_m.       \tag{1}
\]

Consider first $a=3$, so $b=m-3$ is odd and $b\ge5$.
The kernel of the projection of (1) onto $S_b$ is $A_3=C_3$.
Choose a minimal distinguished family generating $S_b$.
The normalization inequality and rank $b$ force its size
$b-1$. Its images are an extremal independent involutory
string. The acquired theorem makes them a natural Coxeter
path, for $b\ge7$; its complete degree-five boundary gives
the same conclusion for $b=5$. They occupy an interval of
the string. The one remaining mark commutes with the whole
path except possibly an endpoint, so its $S_b$ image
centralizes the natural $S_{b-1}$ fixing an endpoint.
That centralizer in $S_b$ is trivial. The remaining involution
would therefore lie in $C_3$, impossible. This excludes $a=3$.

For $b\ge5$ and $a\ne3$, the normal subgroup $1\times A_b$
of $Q$, regarded as a subgroup of the faithful lift $T$,
fixes all $2a$ points over the first orbit. Its action there
maps the perfect group $A_b$ to the abelian group of pair
flips and is therefore trivial. It has no fixed point on
the other $2b$ points, since it is transitive on their pairs.
Its fixed set is exactly the first $2a$ points. A permutation
centralizing $T$ consequently preserves both of these sets.
On the $2b$ points the block image is $S_b$, so the acquired
pair-centralizer lemma says that every centralizer preserves
the pairs. Its exceptional regular prime-dihedral group
cannot be $S_b$ for $b\ge5$. Such a centralizer induces a
permutation centralizing $S_b$ on the blocks, hence induces
the identity. Its remaining flip vector is constant. Thus
its restriction there is either $1$ or $z_b$, the flip of
all $b$ pairs.

If $a\ge4$, the same pair-centralizer lemma applies on the
$2a$ points, with block image $S_a$, and finishes the proof.
If $a=2$, the kernel $A_b$ fixes these four points, so their
induced group is the faithful $S_2$ interchanging the two
pairs. Its involution is one of the two double transpositions
which interchange those pairs. The centralizer of such a
double transposition in $S_4$ is dihedral of order eight;
its even subgroup is the Klein four group of double
transpositions, which preserves every partition into pairs.
Since $b=m-2$ is even, both possible centralizers $1,z_b$
on the other points are even. An even involution centralizing
$T$ must therefore be even on these four points as well,
and preserves their pairs.

The only remaining case has $b<5$, so $m=8$ and $a=b=4$.
The two projections of $Q$ are $S_4$. Their kernels are
normal subgroups contained in $A_4$. Neither is trivial,
as otherwise $Q\cong S_4$ would have five independent
marks. Thus each is $V_4$ or $A_4$.
We claim that the first kernel fixes pointwise the eight
points over the other block orbit. Its action there is a
homomorphism to $\mathbb F_2^4$, equivariant under $Q$.
For $A_4$ it is trivial, since $A_4^{\mathrm{ab}}=C_3$.
For $V_4$, conjugation by the first projection $S_4$ acts
irreducibly on its two-dimensional binary space, so a
nonzero image would have dimension two. Such an image would
be invariant under the second projection $S_4$ acting by
coordinate permutation on $\mathbb F_2^4$. No such subspace
exists: an invariant subspace containing a vector of weight
two contains the entire three-dimensional even subspace;
a vector of weight one or three generates the full space;
the only remaining nonzero vector is the constant vector.
This proves the claim. The first kernel has no fixed points
over its own four-block orbit, because it is transitive on
those four blocks. Its fixed set is therefore exactly the
other eight points. Centralizers of $T$ preserve both
sets. The pair-centralizer lemma, with block image $S_4$
on each set, now proves preservation of all eight pairs.
\end{proof}

\begin{lemma}\label{eight}
Let $Q=(S_4\wr C_2)\cap A_8$ in its natural eight-point
action. Every permutation centralizing a faithful lift $T$
of $Q$ to a permutation group on eight pairs preserves
those pairs.
\end{lemma}
\begin{proof}
The centralizer of $Q$ on its eight letters is trivial.
Indeed its subgroup $A_4\times A_4$ distinguishes the two
four-point sets by its two normal factors; on either set,
the centralizer of the natural $A_4$ is trivial.
If $T$ has two point orbits, each meets every pair once.
The two row actions are the literal $Q$ action. Its trivial
centralizer makes the unique equivariant row interchange
the one matching the original pairs, as required.

If $T$ is transitive, identify it with $Q$. A point stabilizer
$J$ has index two in the stabilizer $M$ of an original pair.
In the eight-letter action,
\[
              M=(S_3\times S_4)^+,
\]
where the superscript denotes equal signs in the two factors.
Its derived subgroup is $A_3\times A_4$. To see equality,
commutation of the first $A_3$ by a lift of a transposition
gives all that $A_3$; the derived subgroup projects onto
$A_4$ in the second factor, since the second projection is
$S_4$. Removing the first coordinate therefore gives all
of the second $A_4$. The opposite containment follows from
the signs. Thus $M^{\mathrm{ab}}=C_2$ and its unique
index-two subgroup is $J=A_3\times A_4$.
This subgroup has orbits of sizes $1,3,4$ on the eight
letters. Its normalizer in $Q$ fixes the unique singleton,
so $N_Q(J)=M$. The coset description of a transitive
centralizer then places every equivariant image of the
chosen point in the same original pair. Equivariance
preserves all pairs, proving the claim.
\end{proof}

\begin{theorem}
Under the standing hypotheses and explicit smaller-degree
induction, $r\le m-1$.
\end{theorem}
\begin{proof}
Assume $r\ge m$, and write $B=\gen L$.
The images of $L$ are independent in $A_m$, so
$|L|\le m-2$. Put $N=B\cap V$ and $K=H\cap V$.
If $C$ is nonempty, its unique mark is $z$; then $m$
is even. IP gives $z\notin B$, so the module list makes
$N=0$. The faithful projection would give a rank-$|L|$
C-marking of $A_m$. Moreover, if $R$ were nonempty,
$\gen{L,R}$ would strictly contain $B$, with the same
block image, and have a nonzero invariant kernel. For even
$m$ every such subspace contains $z$, contradicting IP
with the distinguished set $C$. Hence $R$ is empty and
$|L|=r-2\ge m-2$. This contradicts the acquired finite
or assumed smaller-degree alternating bound.

Thus $C$ is empty. If $R$ were empty too, then
$r=|L|+1\le m-1$. Therefore $R$ has one member and
\[
                   r=m,\qquad |L|=m-2.               \tag{2}
\]
The possibility $N=0$ again contradicts the alternating
bound for the faithful C-marking of $A_m$. The possibility
$N=E$ gives $B=H$, since they have the same block image
and $K\le E$, contrary to independence of the extra mark.
The module list consequently forces
\[
             N=\gen z,qquad K=E,qquad m\text{ even}.
                                                               \tag{3}
\]
The last assertion follows because $z$ must be even.
The interval $L$, the one neighbour occupied by the omitted
label, and the other neighbour occupied by $R$ exhaust the
whole string. After reversal,
\[
 i=0,\qquad L=\{\rho_1,\ldots,\rho_{m-2}\},
 \qquad R=\{\rho_{m-1}\}.
\]
The full parent diagram is connected: otherwise its commuting
components and IP would give a nontrivial direct-product
decomposition of the simple group $A_{2m}$. Every consecutive
pair of marks in the interval therefore fails to commute.

The constant flip $z$ centralizes $H$. The omitted $\rho_0$
does not commute with $z$, since otherwise $z$ would
centralize all of $A_{2m}$, whose natural centralizer is
trivial. Put
\[
                  T=\gen{\rho_2,\ldots,\rho_{m-2}}.
\]
This group centralizes $\rho_0$. If it contained $z$, the
preceding noncommutation would be impossible. By (3) its
block projection is therefore faithful. Denote its image
by $Q\le A_m$. It has the inherited full C-marking of rank
$m-3$, with connected diagram. No quotient intersection
property is assumed: the projection is an isomorphism on
this actual subgroup.

If $m\ge10$ and $Q$ is transitive, the complete L047 theorem
and $m-3>m/2+1$ force $Q=A_m$ or $S_m$. Evenness excludes
$S_m$; the assumed or finite alternating rank bound excludes
$A_m$. If $m=8$, the acquired complete maxima exclude
$A_8$, and the sole proper transitive rank-five exception
in L047 is the order-$576$ action of Lemma~\ref{eight}.
This identification uses its actual accepted permutation
marking: its five generators preserve
\[
             \{1,2,3,8\}\ \big|\ \{4,5,6,7\}
\]
and are even. The even subgroup of that wreath stabilizer
has order $576$, the accepted exact order of the generated
group, so they are equal. This uses the complete retained
census and its actual coordinates, not an abstract group
name alone. Lemma~\ref{eight} makes $\rho_0$ preserve the
original pairs, contradicting primitivity of the parent.

In every remaining case $Q$ is intransitive. Apply
Lemma~\ref{parityproduct} to its faithful lift $T$ and the
even involution $\rho_0$. It again preserves the original
pairs. The group $H$ also preserves them, so
$\gen{H,\rho_0}=A_{2m}$ would be imprimitive. This final
contradiction proves the theorem.
\end{proof}

\paragraph{Attribution and boundary.}
This recovers the alternating block-image case of source Section~3.2.
The source supplies the invariant-module reduction and the rank target;
the new argument passes to the actual faithful endpoint subgroup and
uses normal-factor fixed sets and point-stabilizer normalizers.
It does not assert a splitting of the central extension $B$ and never
regards a block-breaking dihedral group as acting on the blocks.
The complete acquired L047, symmetric extremal and small alternating
classifications retain their mathematical and software premises.
The full source, the symmetric block-image case, and the other general
imprimitive cases remain separate obligations. This component adds no
whole-original, open-question, Lean or compression credit.
\end{document}
